\section{Conclusions}\label{sec:conclusions}
\begin{enumerate}
\item On \NEntries{} maps released in 2026, MICA built an accurate model (median TM-score \TmMedMica{}) for \NModelMica{} entries and did not finish on \NFailedMica{}: two runs were killed by the system for lack of RAM and one reached its time limit, all on a free 31~GB machine.
\item Where both methods built a model (differences are method minus MICA), EModelX(+AF) did not differ from MICA with the same AF3 predictions (\DiffTmEmodelx{}), EMProt was slightly higher (\DiffTmEmprot{}; Holm $p$ \PwHolmTmEmprot{}), E3-CryoFold was lower (\DiffTmEthree{}; Holm $p$ = \PwHolmTmEthree{} for the Wilcoxon test, \PtHolmTmEthree{} for the $t$-test), and ModelAngelo was lower on the maps above 3.0~\AA{} (\DiffTmAngeloHigh{}; Holm $p$ = \PwHolmTmAngeloHigh{} and \PtHolmTmAngeloHigh{}) but not up to 3.0~\AA{} (\DiffTmAngeloLow{}); CryoAtom, added after the first version, did not differ (\DiffTmCryoatom{}; Holm $p$ = \PwHolmTmCryoatom{}), and CryoZeta, added after the first version and run on \NRanCryozeta{} entries, was higher on the \NPairsTmCryozeta{} entries where both built a model (\DiffTmCryozeta{}; Holm $p$ = \PwHolmTmCryozeta{} for the Wilcoxon test, \PtHolmTmCryozeta{} for the $t$-test). Only the differences to EMProt and to CryoZeta are clear on both tests.
\item Counting every failure as zero, the mean TM-score over all \NEntries{} entries is \ZeroMeanMica{} for MICA, \ZeroMeanAngelo{} for ModelAngelo, \ZeroMeanEmodelx{} for EModelX(+AF), \ZeroMeanEmprot{} for EMProt, \ZeroMeanEthree{} for E3-CryoFold and \ZeroMeanCryoatom{} for CryoAtom (for CryoZeta \ZeroMeanCryozeta{} over the \ZeroNCryozeta{} entries on which it ran); MICA does not differ from ModelAngelo, EModelX(+AF), CryoAtom or CryoZeta on this footing, is below EMProt and above E3-CryoFold.
\item CryoZeta, which asks for a GPU of 32~GB, built accurate models where it ran (median TM-score \TmMedCryozeta{}, lowest of the \NModelCryozeta{} models \TmMinCryozeta{}) but finished only \NModelCryozeta{} of the \NEntries{} entries on the free 15~GB machine (\NFailedCryozeta{} failures, \NNotRunCryozeta{} not run); whether it keeps its lead over MICA on the entries where it failed or was not run, and on a larger GPU, is the open question of this part of the benchmark.
\item The step of MICA that uses the AF3 structures raises the median TM-score from \MicaVarInitTmMed{} (the network's own C$\alpha$ model) to \MicaVarAllTmMed{}. Docked input helped slightly on the \DockEffectN{} entries that could be tested (mean difference \DockEffectMean{}).
\item The results hold for the pipeline as I could run it: coarse docking, no real-space refinement, a one-window memory patch, one run per entry. What would tighten them: MICA's own docking on a large-memory machine, its refinement with a Phenix version that accepts its parameter file, a check of the docked placements against the deposited models, repeated AF3 seeds and larger entries.
\end{enumerate}
